\thispagestyle{empty}
\vspace*{0.2cm}

\begin{center}
    \textbf{Zusammenfassung}
\end{center}

\vspace*{0.2cm}

\noindent 
Diese Arbeit zeigt, wie ZoKrates um weitere eingebettete elliptische Kurven erweitert werden kann. Sie gibt eine Einführung in die Arithmetik  von elliptischen Kurven, um sich mit der erforderlichen Terminologie vertraut zu machen.
Dann wird ein SageMath code vorgestellt um eingebettete Kurven zu finden und um mögliche Kandidaten anhand der SafeCurve-Kriterien zu bewerten. Es wird eine detaillierte Implementierung einer eingebetteten Kurve der BLS12\_377s namens Decaf377, sowie eine Signaturüberprüfung in Form einer EdDSA Signatur in der Decaf377 vorgestellt, die innerhalb der BLS12\_377 verschachtelt und wiederum in eine BW6\_761 verschachtelt ist. Anschließend präsentiere ich meine Ergebnisse der SafeCurve Kriterien von bereits implementierten und möglichen zukünftigen Kurven für ZoKrates und schließe die mögliche Implementierung der Bandersnatch, einer eingebetteten Kurve der BLS12\_381, aus. Ich werde einen Überblick über eingebettete elliptische Kurven im Kontext von zk-SNARKs geben und damit die Implementierung neuer eingebetteten elliptischen Kurven für Nicht-Krypto-Experten erleichtern.